\documentclass[11pt]{article}
\usepackage[a4paper,margin=26mm]{geometry}
\usepackage{amsmath,amssymb,amsthm,booktabs,longtable,hyperref}
\newtheorem{theorem}{Theorem}
\newtheorem{lemma}[theorem]{Lemma}
\newtheorem{proposition}[theorem]{Proposition}
\theoremstyle{remark}\newtheorem{remark}[theorem]{Remark}
\newcommand{\Q}{\mathbb Q}\newcommand{\Z}{\mathbb Z}
\newcommand{\Tr}{\operatorname{Tr}}\newcommand{\Nm}{\operatorname{Nm}}
\newcommand{\ord}{\operatorname{ord}}
\title{A scalar-point sharing construction in 163 and 179 variables}
\author{Z diophantine in Q project}
\date{20 November 2026}
\begin{document}\maketitle
\begin{abstract}
We reduce the retained general quartic constructions from 165 to 163
and from 180 to 179 total rational coordinates, including the input.
A partial-trace scalar detector guards its own nonzero quotient.
Scaling the scalar inverse-Vieta chart by that quotient removes its
two slope coordinates. A positive variable weight removes the need
for an additional reciprocal in the general family; the fixed member
uses one reciprocal and retains its effective unweighted height bound. A changed cleared
comparison yields the divisibility $M_j^3\mid a\Delta$, sufficient
for the original global height argument. We give the ordered integer
construction, including its scalar denominators and zero fibers, and
independent exact checks of all rows and expanded coefficients.
The defining assertions retain the explicitly stated arithmetic
interfaces of the source constructions. The 163-variable threshold
is ineffective; the 179-variable coefficient table is fully specified.
\end{abstract}

\paragraph{Source credit.}
The underlying supplied constructions and notes are by Gabe Conant
with Astra~6.2, as audited in~\cite{Shared}. The scalar and point
changes proved here are subsequent project work based on those notes.
The published source tables are retained unchanged.

\section{Statements and retained arithmetic inputs}
Let $K=\Q(\sqrt2)$, put $d_0=3-\sqrt2$, $D=d_0^2=11-6\sqrt2$,
and consider
\[
 E:y^2=x^3-Dx,\qquad R_0=(1+\sqrt2,2\sqrt2),\qquad P_0=2R_0.
\]
For a finite nonzero abscissa write $f(Q)=1/x(Q)$; where the
ordinate is nonzero, write $z(Q)=-x(Q)/y(Q)$.
For the reference point $H$ put
\[
 f=f(H),\quad S=y(H)/x(H),\quad T_H=y(H-P_0)/x(H-P_0).
\]
The scalar normalizer will be denoted by $R$, to distinguish it
from the fixed elliptic point $R_0$.

The arithmetic interfaces are precisely those retained by the
shared-product audit~\cite{Shared}, with the source conventions of
\cite{General,Effective}. More explicitly, we retain:
\begin{enumerate}
\item the group description $E(K)=\Z R_0+E[2](K)$, the complete
finite point charts and label conventions, and the relative
order-four formal-logarithm estimate at consumed good primes;
\item the fixed reference exclusions at 149 and the scalar
contact-height bound, with the same exceptional sets;
\item the full-domain reflected ten-contact height estimates
(center $1$ in the general construction, center $3/22$ in the
effective construction), including the midpoint of radical one;
\item the positive four-point trace-height estimates for every finite
difference allowed by the charts, including finite two-torsion
differences: the height theorem for varying positive weights in
\cite{Weighted} for the general family, and the retained effective
unweighted estimate for the fixed member;
\item the native pole obstruction against negative odd valuations
at consumed primes and its prepared-realization interface for
actual positive inputs below 17, with square reduced denominator
and the prescribed local squares of depths $4,2,2$ at $2,3,17$;
\item the fixed-data Gaussian prime-pattern and compatible
annihilation supplier with arbitrary finite avoidance, arbitrarily
large fixed prime and slope cutoffs, the stated synchronization,
and the finite-index compact-component recurrence used in the
ordered source proof.
\end{enumerate}
The effective assertion additionally retains the numerical
contact-height estimates of~\cite{Effective}. The new scalar
preparation is proved below from the same formal-group and
recurrence inputs; it is not assumed as a replacement interface.
The native realization hypothesis is not removed or inferred from
the finite coefficient calculations.

Set
\[
 \beta=51X^2+2504,\qquad
 r=\beta(\beta-1)(\beta-2401)(32\beta+343).
\]
For the family take
\begin{equation}\label{eq:Rfamily}
 C=149^N,\qquad R=\frac{149C^2r}{1+149^2r},\qquad N\ge1.
\end{equation}
For the fixed member use
\begin{equation}\label{eq:Rfixed}
 R=\frac{A_{\rm a}r}{1+25r},\qquad A_{\rm a}=A_0B_{\rm a}^4,
\end{equation}
where
\begin{align*}
 A_0={}&23\cdot149\cdot263\cdot337\cdot631\cdot1063\cdot1093
 \cdot1171\\
 &{}\cdot1439\cdot2161\cdot2437\cdot2837\cdot2851\cdot3187
 \cdot3343\cdot3793\\
 ={}&645904574641049503754777843805872291231896003979.
\end{align*}
The additional fixed factor is
\[
 B_{\rm a}=263\cdot2161\cdot2437\cdot2837\cdot3793
          =14904184670611231.
\]
Both normalizers are positive on every rational input.

\begin{theorem}\label{thm:main}
Under the respective retained arithmetic interfaces, the exact
tables accompanying this paper define polynomials with rational
input projection $\Z$ as follows.
\begin{enumerate}
\item The family has 163 total variables and defines $\Z$ for
every sufficiently large fixed $N$. Its sufficient threshold is
ineffective. For every $N\ge1$ it has degree exactly four,
137 quadratic residuals, 2789 monomials, content one, and exact
coefficient height
\[
 \max\{126616122944180358034867798936,149^{4N+2}\}.
\]
\item The fixed member has 179 total variables, 151 quadratic
residuals, 3036 monomials, degree exactly four and content one.
Its exact coefficient height is $A_{\rm a}^2$, which has 748 bits.
\end{enumerate}
In both cases the input is coordinate one and the polynomial is
the sum of the squares of the listed ordinary quadratic residuals.
\end{theorem}

The savings concern variable count. The fixed coefficient height
is larger than that of the previous 180-variable member.
The square register $R^2$ remains in the construction; its removal
is not asserted.

\section{The partial-trace scalar detector}
Let $Z\in K$ and define
\begin{equation}\label{eq:scalar}
 J=\frac{Z^2}{R^5}+\frac{f^2}{R^4},\qquad
 \tau\Tr\left(\frac{Z^2}{R^5}\right)=\Nm(J),\qquad
 b_s=R(R+9792)\tau^2.
\end{equation}
Here $f\ne0$ is supplied by the retained reference chart.
At both real embeddings $J>0$. If $Z=0$, the middle equation
would give $0=\Nm(f^2/R^4)>0$. Hence $Z\ne0$ and
\[
 \Tr(Z^2/R^5)>0,\qquad \tau>0,\qquad b_s>0.
\]
In particular the trace denominator is guarded by the equations,
without a new reciprocal coordinate.

\begin{lemma}\label{lem:scalar}
Let $p\nmid2\cdot9792$ and let $\mu=v_p(R)>0$ be odd. If the
native pole obstruction holds for $b_s$, some valuation slot of
$K$ above $p$ satisfies
\[
             v(Z)>\frac{9\mu}{4},\qquad
             v(f)>\frac{7\mu}{4}.
\]
\end{lemma}
\begin{proof}
Normalize $v(p)=1$ and count the two slots with multiplicity at
an inert prime. They are unramified. Write
\[
 a_i=2v_i(Z)-5\mu,\quad b_i=2v_i(f)-4\mu,\quad
 w_i=v_i(J)=\min(a_i,b_i).
\]
The first valuation is odd and the second even, so they cannot
cancel. Since $v_p(b_s)=\mu+2v_p(\tau)$ is odd, the obstruction
gives $v_p(\tau)>-\mu/2$. The trace inequality gives
\begin{align*}
 v_p(\tau)
 &=w_1+w_2-v_p\Tr(Z^2/R^5)\\
 &\le w_1+w_2-\min(a_1,a_2)
 \le w_1+w_2-\min(w_1,w_2)=\max(w_1,w_2).
\end{align*}
Thus both summands in $J$ have valuation greater than $-\mu/2$
at one common slot, proving the inequalities.
\end{proof}

At 149 the known reference norm computations exclude the
forbidden reference strata before any quotient is used. Their
nonzero norms $41^2/64$, $17^2/64$, $63^2/64$ and $7$ are
149-adic units, whereas Lemma~\ref{lem:scalar} forces positive
depth. Hence $S$ and $T_H$ are nonzero. Once the scalar point
is recovered below, $Z=2z(Q_2)/z(H)$ and the retained formal
isometry gives
\begin{equation}\label{eq:scalarindex}
                         v_p(a)>9\mu/4,
 \qquad a=2\nu(Q_2)\in\Z\setminus\{0\}.
\end{equation}
Indeed $H$ is formal at the successful slot. The point $Q_2$
is formal or reduces to $(0,0)$; translation by that two-torsion
point changes $z$ by a sign. Multiplication by two removes the
torsion label, and $h=\nu(H)\in\frac12\Z$ is $p$-integral.
The argument is the source scalar extraction with the same
common-slot bounds, so the scalar radical-height constants are
unchanged.

Put
\[
 \widehat J=R^4J,\qquad \widehat\tau=R^4\tau,
 \qquad L=\widehat J-f^2.
\]
The scalar field equations become
\begin{equation}\label{eq:scalarquadratic}
                 RL=Z^2,\qquad
                 \widehat\tau\Tr(L)=\Nm(\widehat J).
\end{equation}
The field square $f^2$ is already stored by the source detector.
Consequently these are ordinary quadratic component equations.

\section{The point chart and its exceptional fibers}
Use six field coordinates $U,V,W,r_0^*,P^*,r_2^*$, the retained
field ratio $A=L_2$ and reciprocal squares $F_U=f(U_2)^2$,
$F_V=f(V_2)^2$. Define the field expression
\[
                         M=-2S-W
\]
without adding coordinates. Impose
\begin{equation}\label{eq:point}
\begin{aligned}
 UV&=LP^*,&P^*W&=Rr_0^*,\\
 r_0^*(U+V-W-4S)+DRL&=0,&
 LP^*+W(U+V)&=M^2,\\
 r_2^*&=(r_0^*)^2,&F_Ur_2^*&=U^2,\\
 F_Vr_2^*&=V^2,&A(U+M)&=2SZ.
\end{aligned}
\end{equation}
Each displayed row is quadratic over $K$. The rational components
are the rows recorded in the complete physical tables.

\begin{proposition}\label{prop:point}
Equations~\eqref{eq:scalarquadratic} and~\eqref{eq:point} recover
the complete normalized scalar inverse-Vieta chart and its actual
primary and scalar quotients. They exclude none of its required
finite tangent or alternate two-torsion witnesses.
\end{proposition}
\begin{proof}
We already know $R,Z,L\ne0$. If $r_0^*=0$, the sum row gives
$DRL=0$, a contradiction. Thus $P^*,W$ are nonzero by the
product row, and $U,V$ are nonzero by $UV=LP^*$. Define
\[
 (u,v,w,m,\rho)=(U,V,W,M,r_0^*)/Z,
 \qquad p=P^*/R,\quad \rho_2=r_2^*/Z^2.
\]
Direct division gives
\begin{equation}\label{eq:oldpoint}
\begin{aligned}
 uv&=p,&pw&=\rho,\\
 \rho(u+v+w+2m)+D&=0,&p+w(u+v)&=m^2,\\
 \rho_2&=\rho^2,&F_U\rho_2&=u^2,&F_V\rho_2&=v^2.
\end{aligned}
\end{equation}
These are exactly the complete normalized source identities.
They recover the two finite numerator points and their finite
difference $Q_2$, with $p=x(Q_2)$. The last row in
\eqref{eq:point} and the definition of $M$ give
\[
                     A(u+m)=2S,\qquad Z(m+w)=-2S.
\]
Since $S\ne0$, the primary denominator $u+m$ and the scalar
denominator $m+w$ are nonzero. The original point interpretation
therefore recovers $Z=2z(Q_2)/z(H)$. In particular $Q_2$ has
nonzero ordinate, so it is nontorsion under the retained group
description and $a\ne0$.

There is no division by $m+v$: it may vanish, giving the allowed
finite alternate two-torsion point. The complete source point
interpretation also includes tangent secants. Conversely a source
finite nontorsion scalar triple with its nonzero actual $Z$ gives
\eqref{eq:point} by multiplying $u,v,w,m,\rho$ by $Z$, putting
$P^*=Rx(Q_2)$ and $r_2^*=Z^2\rho_2$. The identities above show
that both eliminated components of $M$ then equal $-2S-W$.
\end{proof}

For the general family use the weighted aggregate
\begin{equation}\label{eq:weightedaggregate}
 T=7+\frac14\Tr\bigl(x(Q_1)+R x(Q_2)+x(Q_3)+x(Q_4)\bigr).
\end{equation}
Its scalar-point contribution is $P^*$, so its row is quadratic
without another inverse coordinate. The other three contributions
and all four actual difference points are unchanged. This removes
two rational coordinates in the general family. The uniform weighted
height argument is given below.

The fixed member retains its original unweighted aggregate and the
effective height constants. There $x(Q_2)=P^*/R$, and we add one
rational coordinate $R^{-1}$, guarded by $R R^{-1}=1$, so its two
aggregate components remain quadratic. It therefore saves one
rational coordinate. The variable-weight theorem is not used to
claim an effective 178-variable defining member.

\section{The homogeneous native pole}
Retain the rational register $R_2=R^2$. Set
\[
 \sigma=1/\tau,\quad B=b_s/\tau,\quad
 (E,U_s,V_s,P_s,Q_s)=(e,u_s,v_s,p_s,q_s)/\tau,
 \quad P_2=p_s^2/\tau,\quad Q_2^*=q_s^2/\tau.
\]
The complete quadratic block is
\begin{equation}\label{eq:native}
\begin{aligned}
 \widehat\tau\sigma&=R_2^2,&B\sigma&=R_2+9792R,\\
 \sigma P_2&=P_s^2,&\sigma Q_2^*&=Q_s^2,\\
 3E^2-34E\sigma-289\sigma^2+17BP_2&=0,\\
 E(17\sigma-B)-17U_sV_s&=0,\\
 U_s^2+34U_sV_s-867V_s^2-17BQ_2^*&=0.
\end{aligned}
\end{equation}
The first row guards $\widehat\tau$ and $\sigma$, since $R\ne0$.
With $\tau=\widehat\tau/R^4$, recover
\[
 b_s=\tau B,\qquad
 (e,u_s,v_s,p_s,q_s)=\tau(E,U_s,V_s,P_s,Q_s).
\]
The last three residuals are exactly $\sigma^2$ times
\[
\begin{aligned}
 3e^2-34e-289+17b_sp_s^2&=0,\\
 e(17-b_s)-17u_sv_s&=0,\\
 u_s^2+34u_sv_s-867v_s^2-17b_sq_s^2&=0.
\end{aligned}
\]
These are the retained native pole equations. No native auxiliary
coordinate is inverted, so all of their zero-coordinate branches
are preserved.

\section{The changed comparison and the weighted integer error}
The old cleared comparison would become cubic after the point
scaling. Replace it by the quadratic field expression
\begin{equation}\label{eq:Enew}
       E_2^*=A\overline{(V+M)}+2\overline{T_H}\,\overline Z.
\end{equation}
It uses the existing two error coordinates and their square
registers. At a consumed slot let
\[
 s_2=y(U_2)/x(U_2)=u+m,\qquad
 t_2=y(V_2)/x(V_2)=-(v+m),\qquad
 C_2=A-\overline{2T_H/t_2}.
\]
The reciprocal tests make the indicated points formal and the
local quotients defined. Since $V+M=-Zt_2$, the exact identity is
\begin{equation}\label{eq:Eidentity}
                  E_2^*=\overline{(V+M)}C_2
                         =-\overline{Zt_2}C_2.
\end{equation}

Let $m_2=2\nu(U_2)$, $n_2=2\nu(V_2)$, and $a=m_2-n_2$.
All of these are $p$-integral. The primary reference label is
$h=\nu(H)$ and the opposite label is $h-1$.
At the opposite slot $H-P_0$ is formal. Independently doubling
$R_0$ gives
\[
 z(P_0)=\frac{232+86\sqrt2}{357},\qquad
 \Nm(z(P_0))=\frac{328}{1071}.
\]
Its support is contained in $\{2,3,7,17,41\}$, already excepted.
Hence $z(H)$ is a unit at that opposite slot.

Suppose $v(E_2^*)\ge k$, with $k\le4$. If the opposite valuation
of $z(Q_2)$ is nonpositive, then $V_2$ is finite formal and
\[
 v\bigl(\overline{Zt_2}\bigr)
 =v_{\mathrm{opp}}(z(Q_2))-v_{\mathrm{opp}}(z(H))
                         -v_{\mathrm{opp}}(z(V_2))\le-1.
\]
Thus $v(aC_2)\ge k$. If $z(Q_2)$ has positive opposite
valuation, translate by $(0,0)$ if necessary. The formal
isometry, after multiplication by two, gives
\[
 v_{\mathrm{opp}}(z(Q_2))-v_{\mathrm{opp}}(z(V_2))
                         =v_p(a)-v_p(n_2).
\]
The index $n_2$ is nonzero because $V_2$ is finite formal.
Equation~\eqref{eq:Eidentity} now gives
$v(aC_2)=v(E_2^*)+v_p(n_2)\ge k$ again.

Put $A'=\overline{2T_H/t_2}$. The relative order-four estimates,
used without assuming either ratio integral, state
\[
       hA-m_2=\varepsilon_1,\qquad
       (h-1)A'-n_2=\varepsilon_2,
       \qquad v(\varepsilon_i)\ge4.
\]
Subtracting, and using $C_2=A-A'$, yields
\begin{equation}\label{eq:weighted}
 a(A-a)=-a(h-1)C_2+a(\varepsilon_1-\varepsilon_2),
 \qquad v\bigl(a(A-a)\bigr)\ge k.
\end{equation}
This step precedes integrality transfer.

Use the same first, second and fourth field factors as in the
enhanced source. In the third factor retain the $c^5E_3^4$ term
but replace $-c^7E_2^4$ by $-c^{11}(E_2^*)^4$. Its hyperbolic
pair is now
\[
 \mathcal X=E_3^2+c^3(E_2^*)^2,\qquad
 \mathcal Y=E_3^2-c^3(E_2^*)^2,
 \qquad c^5\mathcal X\mathcal Y
       =c^5E_3^4-c^{11}(E_2^*)^4.
\]
The inverse uses only the already guarded $c\ne0$. All retained
square rows receive this same inverse substitution.
The valuation classes modulo four are unchanged. The existing
shared gate therefore gives a common successful slot where every
factor term has valuation greater than $-\mu/4$. In particular
\[
                    v(E_2^*)>43\mu/16,\qquad v(E_2^*)\ge3.
\]
All other enhanced comparison and homogeneous Lucas precisions
remain four; the input precision remains three.

Let
\[
 \Delta=t^2-(a+1)(a-1)tb+(a-1)^2b^2-1\in\Z
\]
be the source integer Lucas residual. If $v_p(a)\ge3$,
$p^3\mid a\Delta$ is automatic. Otherwise
\eqref{eq:weighted} makes $A-a$ integral of positive valuation,
so $A$ is integral. The other label comparisons transfer their
integrality as in the source. The elementary polynomial difference
between the actual homogeneous Lucas ratio and $\Delta$, using
\eqref{eq:weighted}, gives $p^3\mid a\Delta$. Taking the
squarefree products separately proves
\begin{equation}\label{eq:radicals}
 M_j^3\mid a\Delta\quad(\Delta\ne0),\qquad
 M_j^3\mid J_X\quad(J_X\ne0),
\end{equation}
where $J_X=|\operatorname{num}(\beta-n)|$ is unchanged.
No coprimality between the distinct $M_j$ is asserted. Radical
one, global alternate zero ordinates and other allowed finite
boundary cases are included: the implication at each consumed
prime is proved before its quotients are used.

\section{Global soundness}
All four actual finite difference points and their labels have
been preserved. Put
\[
 B_0=\max(1,|n|,|a|,|t|,|b|).
\]
For the general family apply the positive weighted trace theorem
of~\cite[Theorem 1.1]{Weighted} with the fixed polynomials
$(w_1,w_2,w_3,w_4)=(1,r,1,1)$ and parameter $r=R>0$.
It gives constants $A,J$ such that
\begin{equation}\label{eq:weightedheight}
                  B_0^2\le A\bigl(h(R)+h(T)\bigr)+J.
\end{equation}
The constants depend only on these fixed polynomials, the fixed
curve and the fixed finite-rank group; they do not depend on $N$,
$R$, the input, or any point index.

For clarity, the published family theorem used in that result is
\cite[Theorem 1.1, p.~308]{VB}. Its additive constant is uniform
in the parameter of one fixed projective family and fixed
finite-rank group, but ineffective. Here that family has parameter
$(R,T)\in\mathbf P^1\times\mathbf P^1$ and does not contain $N$
in its defining equations. The proof in~\cite{Weighted} checks
the projectively normal embedding and excludes positive-dimensional
constant translates through every actual tuple: within a conjugate
block the pole coefficients are sums of positive rational multiples
of squares, while between blocks the two conjugate generators are
independent modulo torsion. Torsion coordinates are frozen in
finitely many cases. Thus finite two-torsion points, negative or
repeated indices and the graph boundary create no additional
exception. The required group hypothesis is the one already retained.

Let $\tau_\beta=h(\beta)$. Elementary height inequalities give,
for every rational input and $N\ge1$,
\begin{equation}\label{eq:Rheight}
 h(r)\le4\tau_\beta+\log1801500,\qquad
 h(R)\le4\tau_\beta+(2N+1)\log149+\log1801500.
\end{equation}
Indeed the four linear factors of $r$ contribute row sums
$1,2,2402,375$, and the matrix for
$R=149^{2N+1}r/(1+149^2r)$ has largest row sum
$149^{2N+1}$. The scalar estimates give a fixed $K_0$ with
\[
 L_N=\lfloor9(2N-1)/4\rfloor+1,\quad
 |a|\ge149^{L_N},\quad
 \tau_\beta\le K_0\,149^{-L_N}|a|.
\]
Consequently, uniformly for $B_0\ge|a|\ge149^{L_N}$,
\[
 A h(R)+J= o(B_0^2)\quad\hbox{as }N\longrightarrow\infty.
\]
For example, divide the upper bound from~\eqref{eq:Rheight}
by $B_0^2$: its first term is at most
$4AK_0\,149^{-2L_N}$ and its remaining terms are
$O((N+1)149^{-2L_N})$, with constants fixed before $N$.

If $\Delta\ne0$, then $a\Delta\ne0$ and
$|a\Delta|\le14B_0^5$. The retained two-contact height implication
for any nonzero integer $J_0$ with $M_j^3\mid J_0$, together
with~\eqref{eq:weightedheight}, gives
\[
 B_0^2\le C_1|J_0|^{7/18}+A h(R)+C_2.
\]
Take $J_0=a\Delta$ and first choose $N$ large enough that the
added scalar-height terms are at most $B_0^2/2$. Then
\[
                     B_0^2/2\le C_3B_0^{35/18}.
\]
Since $35/18<2$, $B_0$ is bounded by a constant independent of
$N$. Increasing $N$ now makes $149^{L_N}$ exceed this bound. Thus
$\Delta=0$. The retained Lucas descent gives
\[
 \log B_0\ge(|a|-6)\log(|a|-3),
\]
If $J_X\ne0$, use $M_j\mid J_X$ and
$J_X\le4B_0e^{\tau_\beta}$. The same absorption gives
$B_0^2/2\le C_4B_0^{7/6}e^{7\tau_\beta/6}$, hence
$\log B_0\le(7/5)\tau_\beta+C_5$, with $C_5$ independent of
$N$. This contradicts the Lucas bound for every sufficiently
large fixed $N$. All constants were fixed before this choice;
no height constant depending on $N$ is being hidden in the argument.
Consequently
$\beta=n\in\frac12\Z$. Writing $X=u/v$ in lowest terms gives
$v^2\mid102$, and the squarefreeness of 102 gives $X\in\Z$.

\subsection*{The explicit endpoint}
For every prime $p\mid A_0$, exhaustive exact residue checks give
$r(X)(1+25r(X))\ne0$ for all $X\in\mathbb F_p$. If $v_p(X)<0$,
the leading coefficient $32\cdot51^4$ is a unit and
$v_p(r)=v_p(1+25r)=8v_p(X)$. Thus, for \emph{every rational input},
the fixed member has $v_p(R)=5$ if $p\mid B_{\rm a}$ and
$v_p(R)=1$ at the other anchor primes. Equation~\eqref{eq:scalarindex}
gives $v_p(a)\ge12$ and $v_p(a)\ge3$, respectively. These primes
are disjoint from the scalar radical $M_S(\beta)$, so
\begin{equation}\label{eq:anchorbound}
 \begin{gathered}
 L_{\rm a}M_S(\beta)^3\mid |a|,\qquad
 L_{\rm a}=A_0^3B_{\rm a}^9\ge2^{960},\\
 |a|\ge L_{\rm a},\qquad
 h(\beta)<2^{238}L_{\rm a}^{-1}|a|<|a|/2.
 \end{gathered}
\end{equation}

The effective three-contact constants are retained exactly:
\[
 \mathcal C<2^{300},\quad J_*<2^{75},\quad K_*=2^{304},
 \qquad E(M)<(21/10)M^{1/6006}.
\]
Apply the source reflected-height calculation to $J=a\Delta$.
For $\Delta\ne0$, including $|J|=1$, it gives
\[
 B_0^2<121\mathcal C\sqrt{21/10}\,
 |J|^{1/3+1/36036}
 \left(\log J_*+\tfrac16\log|J|\right).
\]
Using $|J|\le14B_0^5$ and $s=\log B_0$, the exact estimates
\[
 \sqrt{21/10}<3/2,\quad 14^{1/3}<5/2,\quad
 14^{1/36036}<10001/10000,\quad \log2<7/10,\quad\log14<4
\]
give
\begin{equation}\label{eq:endpoint}
       e^{(1/3-5/36036)s}<500\,2^{300}(54+5s/6).
\end{equation}
The third estimate follows already from the first four terms of
the binomial expansion of $(1+1/10000)^{36036}$.
The quotient of the two sides, excluding the constant coefficient,
is increasing for $s\ge0$, because
$1/3-5/36036>5/324$. At $s=960\log2$, its left side divided
by $2^{300}$ is strictly greater than
\[
 2^{20}(1-280/3003)=2^{20}\frac{2723}{3003}>307000,
\]
using $e^{-x}>1-x$, while the right side divided by $2^{300}$
is below $500(54+800\cdot7/10)=307000$. This contradicts
\eqref{eq:endpoint}, so $\Delta=0$.

For the remaining input error, the original weaker $M_j\mid J_X$
suffices. The unchanged coarse bound and~\eqref{eq:anchorbound}
give $s<240+|a|$. The retained Lucas lower bound is larger than
$2(|a|-6)>240+|a|$ for $|a|\ge2^{960}$. Hence $J_X=0$ and
$X\in\Z$ for this actual fixed table.

\section{Ordered integer completeness}
Fix an integer $X$. Every choice below is made before the later
prime data on which it is permitted to depend.

\subsection*{Uniform real neighborhoods}
Let $\mathcal T$ be the identity component of the closure of the
cyclic orbit of $P_0$ in the product of the two real elliptic groups.
Both projections of $\mathcal T$ are nontrivial: a real embedding
cannot make the algebraic nontorsion point $P_0$ torsion. Choose
$H_*$ near the identity in $\mathcal T$, nonzero in both
projections, with $|f(H_*)^{(i)}|<1$, and put $Q_*=H_*$.
The quotient $2z(Q_*)/z(H_*)$ is 2 at both embeddings.
Small fixed neighborhoods therefore give
\begin{equation}\label{eq:realbox}
       |f(H)^{(i)}|\le1,\qquad 1\le|Z^{(i)}|\le3
                         \quad(i=1,2).
\end{equation}
Every finite-index subgroup has the same connected orbit closure.
Include the finite component index in the progressions. Each
resulting residue class then has closure $\mathcal T$, so the
retained recurrence supplies these real returns together with all
finite congruences. After the scalar root is fixed, the finitely
many other chart exclusions can be avoided inside the reference
neighborhood; multiplication by any of its nonzero fixed labels
is nonconstant on $\mathcal T$.

For $R\ge1$,~\eqref{eq:realbox} gives
\[
 J^{(i)}\le\frac{R+9}{R^5},\quad
 \Tr(Z^2/R^5)\ge\frac2{R^5},\quad
 \tau\le\frac{(R+9)^2}{2R^5}.
\]
Consequently
\begin{equation}\label{eq:realbound}
 0<b_s\le\frac{(R+9792)(R+9)^4}{4R^9}
                   \le\frac{24482500}{R^4}<17\qquad(R>100).
\end{equation}
This is a uniform bound independent of the integer input or the
time required for a real return. For~\eqref{eq:Rfamily},
$R\ge149^{2N+1}/22202>100$ when $N\ge1$; for~\eqref{eq:Rfixed},
$R\ge A_{\rm a}/26>100$.

\subsection*{The scalar finite places, including all new denominators}
Put $K_s=R(R+9792)$. For a reduced rational $R=u/v$,
\[
                  K_s=\frac{u(u+9792v)}{v^2},\qquad
                  \gcd(u(u+9792v),v)=1.
\]
Thus every odd valuation of $K_s$ is positive. Prepare the finite
set consisting of all these odd-valuation primes and $2,3,17$.
The source integer congruences give
$v_2(R)\in\{2,3\}$ and $v_3(R)=v_{17}(R)=0$, for both
normalizers. Since $9792=2^6 3^2 17$, the rational $K_s$ is a
square in $\Q_2$, $\Q_3$, and $\Q_{17}$.

First choose $H=[h]P_0$, $64\mid h$, arbitrarily deep in the
formal group at both slots of each prepared prime, while in the
fixed real neighborhood. Include the fixed reduction orders and
component index in this modulus. At each of the finitely many
places choose a positive exact integer depth $t_p$ for
$\lambda=a/h$, large enough for the eventual scalar requirements
but smaller than twice the minimum reference formal depth by an
arbitrarily large fixed margin. Then choose $a=2a_0\in6\Z$ in
these exact residue classes and with $Q_2=[a_0]P_0$ in its real
neighborhood. The recurrence supplies arbitrarily large positive
$a_0$ in each such progression.

Here is the precise local reason that these choices are compatible.
Write $d_w=v_w(z(H))$. On a sufficiently deep formal subgroup,
the formal logarithm has derivative one, $f(H)=z(H)^2$ times a
unit, and
\[
 Z=\lambda(1+\epsilon_w),\qquad
 v_w(\epsilon_w)\ge4d_w-O_p(1).
\]
At good consumed primes this is the retained order-four estimate;
at each fixed preparation prime it also follows by expansion of
the convergent formal logarithm, with a fixed denominator loss.
The scalar root is deeper than $H$, and the rational logarithm
ratio is exactly $a/h$. Given any fixed $q\ge2$, the requirements
\begin{equation}\label{eq:prep}
 v_w(\epsilon_w)\ge q,
 \qquad v_w(Rf^2/\lambda^2)=v_p(R)+4d_w-2t_p\ge q
\end{equation}
hold by taking the reference depth first sufficiently large and
then choosing $t_p$ in the nonempty interval between the fixed
lower requirement and approximately $2\min_w d_w$. Exact depth
is an open congruence condition. Chinese remaindering combines
the finitely many such choices with $3\mid a_0$ and the component
and reduction orders.

Under~\eqref{eq:prep}, $J=\lambda^2R^{-5}$ times a field unit
congruent to one, while
$\Tr(Z^2/R^5)=2\lambda^2R^{-5}$ times a rational unit.
This includes $p=2$, since $q\ge2$ excludes trace cancellation.
Hence the actual input satisfies
\begin{equation}\label{eq:taudepth}
                    v_p(\tau)=2t_p-5v_p(R)-v_p(2).
\end{equation}
Choose the fixed lower requirements for $t_p$ so that $\tau$ is
integral at every odd-support prime and so that $b_s$ has depths
at least $4,2,2$ at $2,3,17$. Its local squareclasses there are
those of $K_s$, so they are the required squares.

No control of new denominator primes of $\tau$ is missing.
Outside the prepared set $v_p(K_s)$ is even, so
\[
                     v_p(b_s)=v_p(K_s)+2v_p(\tau)
\]
is even, however negative $v_p(\tau)$ may be. At an odd-support
prime $v_p(K_s)>0$ and $v_p(\tau)\ge0$, so $v_p(b_s)>0$.
Thus every negative valuation of the actual $b_s$ is even: its
reduced denominator is a square. Together with~\eqref{eq:realbound}
and the local preparations, this establishes every retained native
realization hypothesis for the new scalar input.
After $a_0$ is fixed, all these conditions are open in the reference
point. Shrink to a reference progression preserving them and the
real bounds. This progression is fixed before any comparison
prime is selected.

\subsection*{The remaining roots and the comparison primes}
Choose the Lucas roots and the other three finite point differences
exactly as in the source. With $a\in6\Z$, the retained recurrence
gives positive integral $t$ and odd $b>1$ with $\Delta=0$.
The four points and all shifts are now fixed. In the general family
$Q_1,Q_2$ are nonzero doubles, so at each real embedding their
abscissas exceed $d_0$, while the other two are at least $-d_0$.
Since $R>100$, the weighted sum is greater than $(R-1)d_0>0$,
and~\eqref{eq:weightedaggregate} gives $T>7$, still greater than
the contact diameter four. The fixed member retains exactly the
source shift separation. Enlarge avoidance to include the support of $R$,
the support of the
fixed nonzero coordinates of $Q_2$, the previous preparation
moduli, and the source's fixed denominators and guards.
At every selected prime, $z(Q_2)$ is then a unit at both slots.

Use the same compatible prime supplier and synchronized reference
progression. Its successful primary and opposite formal depths
are equal to some $d\ge1$; only the smaller initial depth is raised.
It is neither necessary nor possible to require both $h$ and
$h-1$ divisible by the prime. At the opposite slot $z(H)$ is a
unit, so~\eqref{eq:Eidentity} has denominator factor of depth
$-d$. The ordinary ratio comparison has depth at least $4d$,
and therefore
\[
                  v(E_2^*)\ge3d,\qquad v(E_3)\ge2d.
\]
Since $v_p(c)=-1$ on these selected witnesses, the changed factor
terms have positive valuations
\[
              v(c^{11}(E_2^*)^4)\ge12d-11>0,
 \qquad v(c^5E_3^4)\ge8d-5>0.
\]
All other terms retain their source positive depths. Consequently
the common product is one modulo $p$ at the successful slot; at
the other slot it is a unit of residue one or has a pole. The
same shared norm/trace gate gives $D_{\rm gate}$ of positive
valuation, and every actual comparison-pole input has selected
valuation one.

At the previously prepared finite places and the real places,
the new test $E_2^*$ is a regular bounded function on the shrunken
neighborhoods. The unchanged contact asymptotics make all weights
tend to zero with pairwise ratios tending to one. The changed
factor still tends to one, as do all other factors, so the actual
comparison inputs preserve their local classes and tend to
$(204/409)^2\in(0,17)$ at the real place. Their parity away from
their own selected supports is unchanged. Apply the retained native
realization theorem independently to the scalar input and to every
comparison input, then fill the complete point charts and all helper
coordinates. The invertible Laurent maps yield rational zeros of
the final quadrics. No earlier root or prime choice is revisited.
This proves integer completeness and finishes Theorem~\ref{thm:main}.

\section{Exact tables, maps and independent verification}
The reproducibility package is
\begin{center}\small
\nolinkurl{computation/variable_reductions_v82/general/}.
\end{center}
Its two complete physical source tables are Laurent polynomials
in $C,c,R$ and the source-labelled ordinary coordinates. The
literal ordinary quadratic tables are in the corresponding
\texttt{quartic163/} and \texttt{quartic179/} directories.
For an ordinary source coordinate the recorded map is
\[
       x_i^{\rm source}=C^{k_i}c^{e_i}R^{r_i}x_i^{\rm stored},
\]
with $C$ omitted in the fixed member. Every row is multiplied by
a nonzero Laurent monomial and a rational scalar to become a
primitive ordinary quadratic. Coordinate one is unchanged.
The common literal guard rows are
\[
 c c^{-1}=1,\qquad c^{-2}=(c^{-1})^2,
 \qquad R_2=R^2.
\]
The fixed member additionally has $R R^{-1}=1$. The family
has no inverse of $R$, so its zero fiber must be checked before
using a Laurent inverse. The input coordinate and the input-square
and scalar pair registers have $R$-weight zero, whereas source
label 101 has $R$-weight $-1$ and $C$-weight one. If $e_{101}$
is its recorded $c$-weight, define
$\rho=c^{e_{101}-1}x_{101}^{\rm stored}$, using only the
already guarded $c$. The ordinary input rows give, even at $R=0$,
\[
 \rho=r(X)>0,\qquad R(1+149^2\rho)=149C^2\rho.
\]
Here the second row is the forced-normalizer row after clearing
its single $R^{-1}$ factor; the other input rows need no power
of $R$ cleared. Therefore $R>0$ follows before any inverse
Laurent substitution. Every recorded monomial map is consequently
invertible on every rational zero, not merely on a generic open
set chosen afterwards.
The hyperbolic inverse is applied in all affected rows.
The two deleted source coordinates are 188 and 189. Source
label 114 stores $R^{-1}$ only in the fixed member.

The defining polynomial is exactly the sum of squares of the
listed rows, so over $\Q$ its zero set is their common zero set.
No zero-set assertion is inferred from numerical sampling.
The expanded coefficient tables, exact support, primitive content
and heights are checked independently. In the family the only
nonconstant coefficient profiles are
\[
                    -298C^2,\qquad-6615898C^2,\qquad22201C^4
\]
on three distinct monomials. For $C=149^N$, $N\ge1$, their
maximum is $149^{4N+2}$; the largest fixed coefficient is the
constant displayed in Theorem~\ref{thm:main}.

From the repository root run
\begin{center}\small
\texttt{python -B computation/variable\_reductions\_v82/general/check.py}.
\end{center}
The default checker writes no files and imports no construction
producer. It reconstructs all retained rows from the literal
published v78 quadratic and weight tables; reconstructs the new
field rows directly over $\Q(\sqrt2)$; reverses every compiled
map; checks the universal point, comparison and native-pole
identities; independently expands both complete sums of squares;
and verifies the exact anchor and endpoint inequalities.
Its immutable receipt pins the dependency and output hashes.
These finite checks certify the stated algebra and arithmetic
inequalities. The global defining theorem follows from the proof
and its explicitly retained arithmetic interfaces.

\begin{thebibliography}{9}\raggedright
\bibitem{Shared}
\emph{Audit and extension of the shared-product quartics}, v78,
\nolinkurl{papers/quartic/shared_product_audit_20261119_v78.tex},
including the enhanced 165/180 tables and their arithmetic interfaces.

\bibitem{General}
\emph{A quartic in 193 variables}, v74,
\nolinkurl{papers/quartic/quartic193_20261117_v74.tex},
especially the complete inverse-Vieta scalar chart, relative formal
errors, unweighted trace height, and ordered synchronized witnesses.

\bibitem{Effective}
\emph{A compact explicit quartic in 229 variables}, v68,
\nolinkurl{papers/quartic/quartic229_compact_20261117_v68.tex},
especially the effective reflected height estimates, scalar radical
bound, three-contact endpoint, and full integer preparation.

\bibitem{Weighted}
\emph{Heights of positive weighted elliptic traces}, v27,
\nolinkurl{papers/trace_heights/weighted_trace_20261114_v27.tex},
Theorem 1.1 and its proof, including variable positive polynomial
weights and all finite torsion strata.

\bibitem{VB}
J. von Buhren, \emph{Borne de hauteur semi-effective pour le
probl\`eme de Mordell--Lang dans une vari\'et\'e ab\'elienne},
Journal de Th\'eorie des Nombres de Bordeaux \textbf{29} (2017),
307--320, Theorem 1.1, printed p.~308.
\url{https://www.numdam.org/article/JTNB_2017__29_1_307_0.pdf}.

\bibitem{LocalCheckpoint}
\emph{A two-term scalar detector and its remaining quadratic-encoding
problem}, v81,
\nolinkurl{computation/variable_reductions_v81/general/README.md}.
The present partial-trace detector replaces that checkpoint's trace
denominator and completes a different point-coordinate saving.
\end{thebibliography}
\end{document}
